Write \(\Phi=\Phi_{J_y,J_z}\), \(G=G_{R,J_y,J_z,n}\), \(H=H^{\mathrm{form}}_{J_y,J_z,n}\), \(B=B_{J_y,J_z,n}\), and \(K=K_{J_y,J_z,n}\).  The tensor wavelets are Lebesgue-orthonormal by \(\mathrm{ass{:}wavelet\text{-}regularity}\).  Hence, for every coefficient vector \(x\),
\[
x^{\top}Gx=\int (x^{\top}\Phi)^2r_n\,dy\,dz,
\qquad c_G\|x\|_2^2\le x^{\top}Gx\le C_G\|x\|_2^2. \tag{1}
\]
For \(d\ge1\), (1) follows from \(\mathrm{ass{:}covariate\text{-}density\text{-}bounds}\); for \(d=0\), it follows from \(r_n\equiv1\).  Since the weight depends only on \(z\), outcome orthogonality gives \(G=I_y\otimes G_z\), so \(G^{-1}\) preserves every outcome level.

For \(g=x^{\top}\Phi\), put \(u_z=\mathcal G_{n,z}g(\cdot,z)\).  Each retained outcome feature has zero outcome integral.  The weak Neumann equation and its boundary conditions give
\[
\int_0^1g(y,z)u_z(y)\,dy
=\int_0^1\bar q_n(y,z)|u_z'(y)|^2\,dy. \tag{2}
\]
For a zero-mean \(v\in H^1(0,1)\), the fundamental theorem of calculus and Cauchy--Schwarz give \(\|v\|_2\le\|v'\|_2\).  The variational characterization of (2), together with \(c\le\bar q_n\le C\) from \(\mathrm{ass{:}conditional\text{-}density\text{-}bounds}\), therefore yields, uniformly in \(z\),
\[
\int_0^1g\,\mathcal G_{n,z}g\,dy\asymp
\|g(\cdot,z)\|_{H^{-1}(0,1)}^2. \tag{3}
\]
The antiderivative of a normalized boundary wavelet at level \(j_y\) has squared \(L^2\)-norm comparable to \(2^{-2j_y}\).  Compact support gives a uniformly bounded overlap number, including the finite boundary block.  Applying (3) level by level and integrating in \(z\) proves
\[
x^{\top}Hx\asymp
\sum_{j_y,k_y,\gamma}2^{-2j_y}x_{j_y,k_y,\gamma}^2. \tag{4}
\]
Equations (1), (4), and preservation of outcome levels by \(G^{-1}\) give
\[
x^{\top}Bx=(G^{-1}x)^{\top}H(G^{-1}x)
\asymp\sum_{j_y,k_y,\gamma}2^{-2j_y}x_{j_y,k_y,\gamma}^2. \tag{5}
\]

For \(f=x^{\top}\Phi\), covariance admits the identity
\[
x^{\top}\Gamma_{a,J_y,J_z,n}x
=\inf_{v\in\mathbb R}\int(f-v)^2q_{a,n}r_n\,dy\,dz. \tag{6}
\]
Because \(\int f(y,z)dy=0\) for every \(z\), Lebesgue orthogonality gives \(\int(f-v)^2dy\,dz=\|x\|_2^2+v^2\).  The density bounds in (6), followed by summation over arms, prove \(c_\Gamma I\preceq\Gamma_{J_y,J_z,n}\preceq C_\Gamma I\).

Let \(w_\lambda=2^{-2j_y(\lambda)}\), ordered decreasingly.  Congruence by \(\Gamma^{1/2}\), (5), and the min--max characterization applied to the spans of the largest reference weights and their orthogonal complements give
\[
\underline c\,w_k^\downarrow\le\nu_{k,J_y,J_z,n}
\le\overline C\,w_k^\downarrow,
\qquad1\le k\le D_{J_y,J_z}. \tag{7}
\]
Taking superlevel sets proves the two counting inequalities.  There are \(\asymp2^{j_y+dJ_z}\) retained indices at outcome level \(j_y\); hence, throughout the asserted threshold range,
\[
N_{J_y,J_z,n}(u)
\asymp2^{dJ_z}\!\sum_{j_0\le j_y\le
\min\{J_y,\frac12\log_2(u^{-1})+O(1)\}}2^{j_y}
\asymp\min\{D_{J_y,J_z},2^{dJ_z}u^{-1/2}\}. \tag{8}
\]
Inverting (7)--(8), with the finite endpoint blocks handled directly, gives
\[
\nu_{k,J_y,J_z,n}\asymp\min\{1,(2^{dJ_z}/k)^2\}.
\]
Consequently
\[
\sum_{k=1}^{D_{J_y,J_z}}\nu_{k,J_y,J_z,n}^2
\asymp2^{dJ_z}+2^{4dJ_z}\sum_{k>2^{dJ_z}}k^{-4}
\asymp2^{dJ_z}. \tag{9}
\]
Equations (7)--(9) prove the operator, Hilbert--Schmidt, and maximal-eigenvalue assertions.  For \(d=0\), the covariate multiplicity in (9) is one, so the Hilbert--Schmidt norm remains bounded above and away from zero.

It remains to retain every cross term.  Let \(P_y\) be the outcome projector and \(P_z^R\) the \(L^2(dR_n)\)-orthogonal covariate projector.  The tensor Gram formula in \(\mathrm{def{:}anisotropic\text{-}projection}\) gives
\[
\Pi_{J_y,J_z}=P_yP_z^R=P_z^RP_y. \tag{10}
\]
Put \(p=P_yP_z^R\Delta_n\), \(r_y=(I-P_y)\Delta_n\), and \(r_z=P_y(I-P_z^R)\Delta_n\).  Thus \(\Delta_n=p+r_z+r_y\).  In the covariate scaling basis, the Gram matrix is banded by compact support and obeys \(cI\preceq G_z\preceq CI\).  With \(\alpha_0=2/(c+C)\) and \(R_z^0=I-\alpha_0G_z\),
\[
\|R_z^0\|_{\mathrm{op}}\le\vartheta:=\frac{C-c}{C+c}<1,
\qquad G_z^{-1}=\alpha_0\sum_{j\ge0}(R_z^0)^j.
\]
Bandedness of \((R_z^0)^j\) and the geometric tail imply
\[
|(G_z^{-1})_{kl}|\le C_1\vartheta^{c_1\operatorname{dist}(k,l)}. \tag{11}
\]
Thus the weighted projector is uniformly local and bounded on the relevant scalar and Banach-valued H\"older classes.  Boundary-wavelet Jackson bounds consequently give
\[
\|r_y\|_2\lesssim2^{-s_yJ_y},\quad
\|r_y\|_{H_y^{-1}}\lesssim2^{-(s_y+1)J_y},\quad
\|r_z\|_2\lesssim2^{-s_zJ_z}, \tag{12}
\]
and \(p,r_z\) have uniformly bounded outcome \(\mathcal C^{s_y}\)-norms.  If \(f\in\{p,r_z\}\), then, writing \(H_f(y,z)=\int_0^yf(v,z)dv\), the Neumann equation gives
\[
\partial_y\mathcal G_nf(y,z)=-H_f(y,z)/\bar q_n(y,z). \tag{13}
\]
The lower density bound, reciprocal H\"older stability, and the extra order furnished by the antiderivative imply
\[
\|(I-P_y)\mathcal G_nf\|_2\lesssim2^{-(s_y+1)J_y},
\qquad
\|(I-P_z^R)P_y\mathcal G_np\|_2\lesssim2^{-s_zJ_z}. \tag{14}
\]

For the symmetric bilinear form \(a(f,g)=\int f\mathcal G_ng\,dy\,dR_n\), self-adjointness and projection orthogonality give
\[
a(p,r_y)=\langle(I-P_y)\mathcal G_np,r_y\rangle,quad
a(r_z,r_y)=\langle(I-P_y)\mathcal G_nr_z,r_y\rangle,
\]
\[
a(p,r_z)=\langle(I-P_z^R)P_y\mathcal G_np,r_z\rangle. \tag{15}
\]
Equations (12)--(15), Cauchy--Schwarz, and (3) give
\[
|a(p,r_y)|+|a(r_z,r_y)|\lesssim2^{-(2s_y+1)J_y},
\qquad |a(p,r_z)|+a(r_z,r_z)\lesssim2^{-2s_zJ_z}, \tag{16}
\]
and \(a(r_y,r_y)\lesssim2^{-2(s_y+1)J_y}\).  Finally,
\[
B^y_{J_y,n}=2a(p,r_y)+2a(r_z,r_y)+a(r_y,r_y),
\qquad
B^z_{J_z,n}=2a(p,r_z)+a(r_z,r_z). \tag{17}
\]
Equations (16)--(17) prove the two claimed envelopes, including the retained--omitted cross terms.  Since (17) compares full and projected energies, it also proves that these quantities are approximation gaps rather than biases of the exact-target estimator.